\documentclass[11pt,a4paper]{article}
\usepackage[utf8]{inputenc}
\usepackage[T1]{fontenc}
\usepackage[margin=1in]{geometry}
\usepackage{fvextra}
\RecustomVerbatimEnvironment{verbatim}{Verbatim}{
  breaklines=true,
  breakanywhere=true,
  fontsize=\small
}

\title{Report - 25.02.2026}
\date{25.02.2026}

\begin{document}
\maketitle

\begin{verbatim}
First rule check ✅: some older uploads from earlier in this chat session are expired, but for what you asked **right now** we’re good — I have `main.tex`, all **final clean CSVs**, and the core PDFs (DESTINO + your consolidated literature notes). If you later want me to use a specific expired PDF again, you’ll need to re-upload that one.

I opened `main.tex` and read it line-by-line, and I cross-checked every number in your tables against the **final clean datasets** you produced today (`*_clean.csv`). Your tables are **actually correct** (good news). Below is the deep evaluation you asked for: roadmap alignment, DESTINO alignment, data quality, loose ends, and improvement plan.

---

# 1) How well did we operate vs the roadmap?

Your professors’ roadmap (“second bible”) is:

1. literature review ✅
2. define comparison criteria incl feasibility ✅
3. select OSS projects ✅
4. compute & compare multiple metrics ✅ (compute done; compare deferred correctly)
5. select one metric for PoC ⛔ (correctly not done yet)
6. PoC smart contract ⛔ (correctly not done yet)
7. discussion/writing 🟡 (Phase I report is solid and now consistent)

**Roadmap compliance score: 9.5/10**

Why not 10? Only because we still have one unresolved extraction limitation (timestamp semantics) that you already documented, and we should either fix it or plan sensitivity analysis.

---

# 2) Is the report aligned with DESTINO?

DESTINO (“primary bible”) is mainly about evidence: **hash + link + signature**, then store **metrics + tags + storage persistence**, and keep things auditable.

### What is aligned ✅

Your report reflects DESTINO-friendly thinking in multiple places:

* You explicitly state **no PoC metric is selected yet** and selection happens after Phase II (good: matches professor constraint and prevents premature bias).
* You consistently frame feasibility as “hybrid is recommended” (off-chain computation + on-chain compact commitments), which is the only realistic way to implement this within constraints.
* You include **storage persistence** and **software-type tags** in the dataset freeze (DESTINO metadata requirement).
* You keep the pipeline **auditable**: frozen repos/tags/windows, cleaning logs, and canonical datasets.

### What is not yet implemented (but that’s OK at Phase I) ⚠️

DESTINO’s on-chain workflow isn’t implemented yet (signature verification, storing hashes on ledger, emitting events). But Phase I is supposed to end before PoC. So this isn’t a failure — it’s just **future work**.

**DESTINO alignment at Phase I score: 8/10**

If you want 9–10/10 later, you’ll add one section in Month 4 describing how your PoC stores:

* `archive_hash`, `archive_link`, `signer_id`, `tags`, `storage_persistence`, plus `metric_value` and `evidence_hash`.

---

# 3) Deep data audit vs the final CSVs

I validated your final clean bundle:

### Dataset shape (matches exactly what we planned)

* 4 repos frozen ✅
* 8 stable tags per repo → **32 tags total** ✅ (`versions.csv`)
* 7 windows per repo → **28 windows total** ✅ (`windows.csv`)
* LOC per tag → **32 rows** ✅ (`loc_per_tag.csv`)
* clean developer-window dataset → **273 rows** ✅ (`metrics_window_dev_clean.csv`)
* clean window totals → **28 rows** ✅ (`metrics_window_totals_clean.csv`)
* clean ownership dev shares → **273 rows** ✅ (`ownership_dev_shares_clean.csv`)
* clean ownership summaries → **28 rows** ✅ (`ownership_summary_clean.csv`)

### Your report tables are correct

I re-computed and matched the numbers you printed, including:

* min/max/mean window durations per repo (e.g., requests min window 0.26 days) ✅
* total commits and churn per repo ✅
* global totals: **902 commits**, **58,298 churn** ✅
* LOC min/max/mean per repo ✅
* M4 concentration means and gini values ✅

So the “tables aren’t correct” fear is understandable (because in past versions they were), but **in this current `main.tex` the numbers match the final clean CSVs**.

**Data correctness score: 9/10**

The remaining 1 point is exactly one known limitation:

### Loose end (still true): timestamp semantics (“author-date drift”)

You correctly document that some `first_commit` values can be earlier than the window boundary because timestamps are author-date, not committer-date. This doesn’t break churn or commit counts, but it can distort cadence in some analyses.

You already wrote it as a known limitation, which is good.

---

# 4) What are the loose ends we must fix or attend to?

These are the real ones (not “nice to have”).

## Loose end A — Decide how you’ll handle cadence timestamps

In your report you plan “sensitivity checks on timestamp semantics” — great.

What you must do later (Phase II):

* either re-extract cadence using committer-date (`%cI`) OR
* keep author-date but treat cadence primarily as “commit counts per window” + “per-day normalized rates” and mention timestamp limitation.

## Loose end B — Normalize by window duration in Phase II (you already noted this)

Windows vary wildly:

* requests has a 0.26-day window (!!)
* click has windows up to ~490 days

So Phase II must use:

* commits/day
* churn/day
* optionally churn/(KLOC·day)

You already mention this in Section `\ref{sec:evaluation_unit}` and in the window duration table; that’s exactly correct. The loose end is: **make sure you actually do it in Month 3**.

## Loose end C — “M5 maintenance responsiveness”

Right now M5 is listed as “design extension” — that’s fine.
But to avoid confusion, tweak wording slightly:

* In shortlist table you say: “not implemented in baseline PoC”
* Better: “not computed in Phase I; documented as PoC extension”
  Because PoC isn’t started yet.

## Loose end D — References file / compilation

Your `main.tex` ends with:

```latex
\nocite{*}
\bibliographystyle{plain}
\bibliography{references}
```

If your `references.bib` is missing or only contains 1 entry, the PDF will show “1 reference.”

So the loose end is not your LaTeX — it’s the BibTeX file and compilation steps.
Make sure the project includes `references.bib` with entries for:

* destino
* blockchainsurvey
* smartmetrics2023
* scvulncomplexity
  …and that you compile: `pdflatex → bibtex → pdflatex → pdflatex`.

---

# 5) Where can the report improve (and why)?

Your report is already a strong **Phase I technical report**. Improvements now are mainly about clarity and “thesis polish.”

## Improvement 1 — Reduce duplicated phrasing, keep structure tighter

Right now you explain “Phase I freeze/pipeline/cleaning” in multiple sections (which is okay), but you can tighten by:

* Keeping one “Phase I summary” paragraph at the end of each chapter
* And referencing the execution ledger table (`\ref{tab:phase1_execution_ledger}`) instead of repeating steps

This improves readability without losing content.

## Improvement 2 — Add a “Data Dictionary” section (high value)

This is a classic thesis improvement: define exactly what each CSV column means.

For example:

* `metrics_window_dev_clean.csv` columns: commits, added, deleted, churn, active_days, is_bot, etc.
* `ownership_summary_clean.csv`: top1/top3 share definitions, gini meaning

This makes the thesis readable for non-experts and is a common examiner expectation.

## Improvement 3 — Make DESTINO mapping explicit (one short subsection)

Add a small subsection like:

* “How Phase I datasets map to DESTINO ledger records”
* Mention: these CSV outputs are what later become “metric records” committed on chain (hash + metadata + value).

That will raise DESTINO alignment score.

## Improvement 4 — Replace “N/A” in rubric template table

In `Table~\ref{tab:rubric_template}`, you put “N/A” in every cell.
That can confuse a reviewer (they might think “not applicable” rather than “not scored yet”).

Better: leave cells blank or put “(deferred)” once in caption.

---

# 6) Scorecard for what we’ve done so far

Here’s a fair grading for Phase I:

* **Roadmap discipline:** 9.5/10
* **DESTINO alignment (phase-appropriate):** 8/10
* **Dataset freeze quality:** 10/10
* **Metric extraction correctness (M1–M4):** 9/10
* **Cleaning/validation rigor:** 9.5/10
* **Report quality as Phase I technical report:** 8.5/10
  (mainly limited by bibliography wiring + missing data dictionary + small wording tweaks)

Overall: **9/10 for Phase I execution**.

---

## Next thing to fix (before you pause)

If you want the fastest “quality jump” without starting Month 3:

1. add a **Data Dictionary** section
2. ensure `references.bib` exists + contains all citations + compiles
3. tweak M5 wording (“not computed in Phase I; design extension”)

If you want, paste your `references.bib` (or upload it), and I’ll verify the bibliography keys match the `\cite{...}` keys in your `main.tex` so you don’t get empty references.
\end{verbatim}

\newpage
\section{Solutions in Details}
\begin{verbatim}
First rule check ✅: **some older uploads have expired in this chat session.** For the loose ends we’re addressing now, we’re still fine because we have your **final clean CSVs**, `main.tex`, `DESTINO4MasterThesis.pdf`, and your consolidated literature notes. If you want me to use/cite any specific expired paper again, just re-upload it.

Now I’ll “zoom in” on each loose end exactly as you asked:

1. why it’s a loose end
2. how it affects the project
3. top 3 best solutions
4. what to do (detailed steps)

---

# Loose end 1 — Timestamp semantics (“author date drift”)

## 1) Why it’s a loose end

In your extracted dataset, some time fields (`first_commit`, `last_commit`, and therefore `active_days`) are based on Git **author date**. In real Git history, author date can be much older than when the commit actually entered the branch (e.g., merges, cherry-picks, rebases). That’s why you observed cases where `first_commit` appears *before* the window’s `from_date`.

This is a loose end because Month 3 includes **cadence/stability over time**, and time semantics must be consistent to avoid misleading results.

## 2) How it affects our project

* **Does NOT affect:** M2 churn, commit counts, window membership (the commit range `from_hash..to_hash` is correct).
* **Can affect:**

  * `active_days` (if author date is old, active days may fall outside the window)
  * any “time gap” and cadence interpretation using `first_commit/last_commit`
  * any “per-day” normalization if you accidentally use timestamps rather than window duration

If we ignore it, the professor could fairly say:

> “Your cadence metric is not well-defined; you’re mixing timestamps that don’t correspond to window boundaries.”

## 3) Top 3 best solutions (with detailed steps)

### Solution A (Best, simplest, most defensible): Re-extract cadence timestamps using **committer date**

**Why it solves it:** committer date reflects when the commit was actually applied to the branch history. It’s the right choice for “activity in this window”.

**What to do (exactly):**

1. Modify the extraction command in your script to output **committer date** instead of author date.

   * Replace `%ad` with `%cI` (ISO 8601 committer date).
2. Recompute for each window:

   * `active_days` from committer date
   * `first_commit` / `last_commit` from committer date
3. Keep everything else the same (added/deleted/churn remains unchanged).
4. Re-run:

   * `extract_window_metrics.py` → produce `metrics_window_dev_committer.csv`
   * then run the same cleaning script → `metrics_window_dev_clean_committer.csv`
5. In the thesis, state:

   * “Cadence features are based on committer timestamps; churn is based on diff stats; windows are tag-to-tag.”

**Pros:** most accurate, clean story, no ambiguity.
**Cons:** you’ll regenerate 2–3 CSVs (but it’s fast).

---

### Solution B (Good, minimal disruption): Keep author-date but redefine cadence to be **count-based only**

**Why it solves it:** you avoid relying on questionable timestamps; cadence becomes “number of commits per window” and “commit rate using window duration”, not timestamp-based first/last.

**What to do:**

1. In Month 3, define cadence as:

   * `commits_per_window`
   * `commits_per_day = commits / window_duration_days` (window duration comes from tag dates, not commit timestamps)
2. Drop or downplay:

   * `first_commit`, `last_commit`, `active_days`
3. In your thesis:

   * explicitly say author-date drift exists; hence cadence is operationalized using counts and window length, not raw commit timestamps.

**Pros:** no re-extraction.
**Cons:** you lose “active days” nuance; cadence becomes simpler.

---

### Solution C (Strongest academically, slightly more work): Dual-run robustness check (author vs committer)

**Why it solves it:** you can show the professor you tested sensitivity. If results are similar, you gain credibility.

**What to do:**

1. Keep the current dataset as “author-date cadence”.
2. Create a second dataset “committer-date cadence” (Solution A extraction).
3. Run Month 3 cadence analyses twice (or at least compare summary stats):

   * stability metrics with author-date and committer-date
4. Report in thesis:

   * “Results are robust (or not) to timestamp semantics.”

**Pros:** best scientific practice; strong defense.
**Cons:** more work and more reporting.

---

# Loose end 2 — Window duration variability (very short vs very long windows)

## 1) Why it’s a loose end

Because your windows are defined by release tags, the time between releases is not constant. Some windows are tiny (hours/days), others are months. That makes raw counts misleading.

Example intuition:

* 50 commits in 2 days ≠ 50 commits in 200 days (very different intensity)

This is a loose end because Month 3 includes cross-window comparisons and “stability across versions.”

## 2) How it affects our project

* Without normalization, metrics like churn/commits will look “unstable” simply because one window is longer.
* Cross-project comparisons become unfair because projects have different release rhythms.
* It also affects ownership concentration: long windows accumulate more contributors.

## 3) Top 3 best solutions

### Solution A (Best default): Rate-normalize by window duration (days)

**Why it solves it:** makes windows comparable.

**What to do:**

1. Compute `duration_days = (to_date - from_date)`.
2. Add derived metrics in Month 3:

   * `commits_per_day`
   * `churn_per_day`
   * optionally `active_days_per_day = active_days / duration_days` (if you fix timestamps)
3. Use rate metrics for stability comparisons (RQ1), while still reporting raw totals as context.

**Pros:** simple, standard, defensible.
**Cons:** very short windows can inflate rates (need guard).

---

### Solution B (Also strong): Use fixed time windows (monthly) instead of tag windows (for stability)

**Why it solves it:** makes all windows equal length.

**What to do:**

1. Keep tag windows for “release-to-release story.”
2. Add a parallel monthly-window dataset (12 months) for cadence stability.
3. Compare stability using monthly windows; use tag windows as secondary.

**Pros:** very comparable and stable.
**Cons:** adds dataset complexity; you already froze tag windows.

---

### Solution C (Minimal): Keep tag windows, but apply **minimum window length rule**

**Why it solves it:** prevents nonsense windows (like <1 day) from distorting results.

**What to do:**

1. If `duration_days < X` (e.g. 7 days):

   * merge it with the next window, or
   * treat as an “outlier window” and exclude from stability stats (documented)
2. Justify: extremely short release gaps produce noisy windows.

**Pros:** simple, keeps tag framing.
**Cons:** changes frozen windows; must be documented carefully.

---

# Loose end 3 — Ownership results “bots” and identity consistency

## 1) Why it’s a loose end

You correctly flagged bots (`is_bot`). But if you don’t explicitly decide how bots are treated in results, you risk confusion:

* Bots can contribute churn/commits (dependency bumps, formatting) which distort ownership and concentration.

Identity consistency (email aliases) was mostly fixed via cleaning, but across windows a person could still have multiple emails (future issue in Round 2).

## 2) How it affects the project

* Ownership concentration (top1/top3, gini) can be inflated if bots produce lots of small commits.
* Professor may ask: “Are you measuring developers or automation?”

## 3) Top 3 best solutions

### Solution A (Best practice): Report both “with bots” and “without bots”

**Why it solves it:** shows robustness and transparency.

**What to do:**

1. In Month 3, compute ownership stats twice:

   * including bots
   * excluding bots (`is_bot=False`)
2. In thesis: show main results (no-bot) and a sensitivity note (with-bot).

**Pros:** strongest, defensible.
**Cons:** slightly more reporting.

---

### Solution B (Minimal): Exclude bots from all ownership calculations (keep in raw data)

**Why it solves it:** keeps ownership about humans.

**What to do:**

1. Filter rows where `is_bot=True` before calculating:

   * shares
   * top1/top3
   * gini
2. Still keep bot rows in raw dataset for auditability.

**Pros:** clean ownership story.
**Cons:** you must justify why bots are excluded.

---

### Solution C (Advanced): Keep bots, but down-weight them

**Why it solves it:** recognizes bots contribute but don’t represent developer quality.

**What to do:**

1. Assign `bot_weight = 0.1` and multiply bot commits/churn before computing shares.
2. Explain as a heuristic; include sensitivity results.

**Pros:** nuanced.
**Cons:** heuristic may be challenged; I don’t recommend for minimum thesis.

---

# Loose end 4 — Bibliography and compilation (why only “1 reference” shows up)

## 1) Why it’s a loose end

Your `main.tex` can be correct but if:

* `references.bib` is missing, or
* BibTeX isn’t run, or
* there are no \cite calls,
  the references section will appear empty or with 1 item.

## 2) How it affects the project

* A thesis with missing references looks unfinished and can hurt evaluation immediately.
* It breaks credibility: claims aren’t anchored.

## 3) Top 3 best solutions

### Solution A (Best): Use BibTeX correctly and add real citations

**What to do:**

1. Ensure `references.bib` contains at least these 4 entries:

   * DESTINO proposal
   * blockchain survey
   * PLOS smart-metrics
   * smart-contract complexity/vulnerability paper
2. Add 6–10 citations in the text.
3. Compile in this order:

   * pdflatex
   * bibtex
   * pdflatex
   * pdflatex

---

### Solution B (Debug mode): add `\nocite{*}` temporarily

**What to do:**

1. Add `\nocite{*}` to force printing all bib entries.
2. Once stable, remove it and keep only cited refs.

---

### Solution C (If BibTeX is painful): switch to `biblatex` + biber

**What to do:**

1. Use `biblatex` package and `\addbibresource{references.bib}`
2. Compile with `biber` instead of `bibtex`

---

# My recommendation (minimum work, maximum thesis strength)

If you want the cleanest path with minimum risk:

1. **Fix cadence timestamps with committer date** (Solution 1A)
2. **Use rate-normalized metrics** for stability (Solution 2A)
3. **Report ownership both with and without bots** (Solution 3A)
4. **Fix bibliography** with BibTeX (Solution 4A)

If you say “go”, I’ll give you **exact Codex prompts** for:

* re-extracting using committer date,
* generating rate-normalized derived columns,
* generating bot/no-bot ownership summaries,
* and verifying BibTeX compilation.
\end{verbatim}

\end{document}


